\uuid{n7Ru}
\titre{Loi de Cauchy et existence de l'espérance}
\auteur{Maxime Nguyen}
\organisation{AMSCC}
\datecreate{2026-11-24}
\niveau{L2}
\module{Probabilité}
\chapitre{Probabilité continue}
\sousChapitre{Lois usuelles et approximations}
\difficulte{2}
\theme{loi de Cauchy, existence de l'espérance, symétrie}

\contenu{
\begin{question}
Soit $X$ de densité $f(x) = \dfrac{1}{\pi(1 + x^2)}$ (loi de Cauchy de centre $x_0 = 0$ et de paramètre d'échelle $\lambda = 1$ ; sa densité est symétrique par rapport à $x_0$). Quelles affirmations sont vraies ?
\begin{answers}
\bad{$f$ est paire, donc $\mathbb{E}(X) = 0$.}
\good{$X$ n'a pas d'espérance, car $\displaystyle\int_0^{+\infty} x f(x)\,\mathrm{d}x = +\infty$.}
\good{$\mathbb{P}(X \leq 0) = \frac12$.}
\bad{$f$ n'est pas une densité, puisque $X$ n'a pas d'espérance.}
\end{answers}
\begin{explanations}
$\int_0^{+\infty} \frac{x}{\pi(1+x^2)}\,\mathrm{d}x = +\infty$,
donc $x \mapsto x f(x)$ n'est pas intégrable. L'argument de parité ne
s'applique qu'à une fonction intégrable. En revanche la symétrie donne bien
C, et $f$ est une densité (positive, d'intégrale 1).

Contre-sens visé : « la densité est symétrique autour de $x_0$, donc l'espérance vaut $x_0$ »,
sans vérifier l'existence de l'espérance.
\end{explanations}
\end{question}
}
